\documentclass[11pt]{article}

\usepackage[T1]{fontenc}
\usepackage[utf8]{inputenc}
\usepackage{lmodern}
\usepackage[a4paper,margin=27mm]{geometry}
\usepackage{microtype}
\usepackage{graphicx}
\usepackage{amsmath,amssymb}
\usepackage[table]{xcolor}
\usepackage{booktabs}
\usepackage{caption}
\usepackage{enumitem}
\usepackage{listings}
\usepackage{fancyhdr}
\usepackage{hyperref}

\definecolor{EvoGreen}{HTML}{168761}
\definecolor{EvoDark}{HTML}{20272C}
\definecolor{EvoMuted}{HTML}{5F6B73}
\definecolor{EvoPanel}{HTML}{F2F6F4}
\definecolor{EvoWarning}{HTML}{FFF8DE}
\definecolor{EvoRule}{HTML}{CBD5D0}

\hypersetup{
  colorlinks=true,
  linkcolor=EvoGreen,
  urlcolor=EvoGreen,
  citecolor=EvoGreen,
  pdftitle={EvoXRB v0.2.0: Watching a Genetic Algorithm Fit an X-ray Spectrum},
  pdfauthor={Sam Seifi},
  pdfsubject={A continuation of the EvoXRB project blog},
  pdfkeywords={genetic algorithm, X-ray astronomy, SciPy, Matplotlib, spectral fitting}
}

\pagestyle{fancy}
\fancyhf{}
\fancyhead[L]{\small\textcolor{EvoMuted}{EvoXRB / Part II}}
\fancyhead[R]{\small\textcolor{EvoMuted}{v0.2.0}}
\fancyfoot[C]{\small\thepage}
\renewcommand{\headrulewidth}{0.4pt}
\renewcommand{\headrule}{\hbox to\headwidth{\color{EvoRule}\leaders\hrule height \headrulewidth\hfill}}
\setlength{\headheight}{14pt}

\setlength{\parindent}{0pt}
\setlength{\parskip}{0.78em}
\setlist{itemsep=0.35em,topsep=0.4em}
\captionsetup{font=small,labelfont=bf}
\graphicspath{{../assets/}}

% Support the documented article-directory build and a build launched from the
% EvoXRB repository root.
\def\EvoComparisonPath{}
\IfFileExists{../assets/e08-ga-comparison.png}{%
  \renewcommand{\EvoComparisonPath}{../assets/e08-ga-comparison.png}%
}{%
  \IfFileExists{website/evoxrb-v0.2.0/assets/e08-ga-comparison.png}{%
    \renewcommand{\EvoComparisonPath}{website/evoxrb-v0.2.0/assets/e08-ga-comparison.png}%
  }{}%
}

\lstdefinestyle{evoxrb}{
  basicstyle=\ttfamily\small,
  backgroundcolor=\color{EvoPanel},
  frame=single,
  rulecolor=\color{EvoRule},
  breaklines=true,
  columns=fullflexible,
  keepspaces=true,
  showstringspaces=false,
  xleftmargin=0.5em,
  xrightmargin=0.5em,
  aboveskip=0.8em,
  belowskip=0.8em
}
\lstset{style=evoxrb}

\newenvironment{scopebox}[1][Scope]
{%
  \begin{center}
  \begin{minipage}{0.92\linewidth}
  \color{EvoDark}
  \hrule height 1pt
  \vspace{0.55em}
  \textbf{\textcolor{EvoGreen}{#1}}\par
  \small
}
{%
  \par\vspace{0.3em}
  \hrule height 1pt
  \end{minipage}
  \end{center}
}

\newenvironment{warningbox}
{%
  \begin{center}
  \begin{minipage}{0.92\linewidth}
  \color{EvoDark}
  \hrule height 1pt
  \vspace{0.55em}
  \textbf{\textcolor{EvoGreen}{Calibration boundary}}\par
  \small
}
{%
  \par\vspace{0.3em}
  \hrule height 1pt
  \end{minipage}
  \end{center}
}

\title{\vspace{-1.6cm}
  {\small\color{EvoGreen}\MakeUppercase{Computational astrophysics / Project blog / Part II}}\\[0.8em]
  \textbf{EvoXRB v0.2.0: Watching a Genetic Algorithm Fit an X-ray Spectrum}
}
\author{Sam Seifi}
\date{August 2026}

\begin{document}
\maketitle

\begin{center}
\begin{minipage}{0.86\linewidth}
\large
The first version saved the winning numbers. The second lets me watch thirty-two
synthetic organisms argue with a Poisson spectrum, one generation at a time.
\end{minipage}
\end{center}

\begin{scopebox}
EvoXRB v0.2.0 remains a declared synthetic / NICER-inspired experiment. The
response, background and target counts discussed here are educational products
generated in Python. The purple open circles are a separate real MAXI/GSC
spectrum of MAXI J1820+070. They provide observational context, but they are
never used in the C-statistic or residuals and are not a real NICER fit.
\end{scopebox}

\section*{Introduction}

The first EvoXRB post ended with a sentence I liked:

\begin{quote}
\textit{You can watch the search happen.}
\end{quote}

There was only one small problem.

You could not actually watch the search happen.

Version 0.1.0 saved convergence histories, gene spreads and final figures. That
was enough to reconstruct the behaviour of the genetic algorithm after the
calculation had finished. It was rather like receiving a very detailed report
from an expedition while being told that the windows on the vehicle were an
optional feature planned for the next release.

Version 0.2.0 adds the windows.

Every evaluated generation can now emit a snapshot of the population. The best
genes in that snapshot are decoded back into astrophysical parameters, passed
through the same photon model, folded through the same synthetic response and
drawn against the same Poisson counts used by the objective.

The result can update live while the optimiser is calculating. It can also be
saved as a replay with play, pause, step, loop and timeline controls.

This sounds like a presentation feature.

It is.

It also became a much more direct test of whether the calculation was doing what
I thought it was doing.

\section{What actually changed in v0.2.0?}

The optimiser remains a real-valued genetic algorithm. There is no prerecorded
sequence pretending to evolve and no linearly interpolated spectrum being
nudged towards the answer because it looks pleasing in a GIF.

The population still passes through tournament selection, simulated-binary
crossover, bounded polynomial mutation, elitism and, when stagnation rules
allow it, random immigration. Every candidate is decoded from a normalized
chromosome

\begin{equation}
  \mathbf{g} = [g_1,g_2,g_3,g_4], \qquad 0 \leq g_j \leq 1,
\end{equation}

into the physical parameter vector

\begin{equation}
  \boldsymbol{\theta} =
  [T_{\mathrm{in}},N_{\mathrm{disk}},\Gamma,K].
\end{equation}

The two normalizations use logarithmic transforms, because treating a mutation
near 1 and a mutation near 10,000 as the same additive event is a good way to
turn numerical conditioning into performance art.

The new part is observational rather than reproductive. The optimiser reports
its state after generation zero and after every newly evaluated generation.
Each snapshot contains:

\begin{itemize}
  \item the normalized population and its scores;
  \item the best genes and their decoded parameters;
  \item the best and median C-statistic;
  \item gene spread and boundary-hit diagnostics;
  \item the evaluation count, random seed and stopping state.
\end{itemize}

Those arrays are independent read-only copies. A callback may inspect the gene
pool, draw it or preserve it. It cannot reach back into the optimiser and give
its favourite organism a slightly more advantageous disk temperature.

That would be less natural selection and more suspiciously interventionist
gardening.

\section{Turning a generation into an X-ray spectrum}

A population history is not automatically a spectral history.

The GA works in normalized coordinates. The detector plot needs expected counts
per channel. Connecting the two requires the full forward model for every frame:

\begin{equation}
  \mathbf{g}^{(m)}_{\mathrm{best}}
  \longrightarrow
  \boldsymbol{\theta}^{(m)}_{\mathrm{best}}
  \longrightarrow
  S(E\mid\boldsymbol{\theta}^{(m)}_{\mathrm{best}})
  \longrightarrow
  \boldsymbol{\mu}^{(m)}.
\end{equation}

Here $m$ is the generation index, $S(E\mid\boldsymbol{\theta})$ is the
educational photon model and $\boldsymbol{\mu}$ is the response-folded expected
detector count vector.

The folding step remains

\begin{equation}
  \mu_i = t_{\mathrm{exp}}
  \sum_j R_{ij}\,A_j\,S(E_j\mid\boldsymbol{\theta})\,\Delta E_j
  + t_{\mathrm{exp}}\,b_i\,\Delta E_i,
\end{equation}

where $A_j$ is the effective area, $R_{ij}$ is the redistribution probability
from true-energy bin $j$ to detector channel $i$, and $b_i$ is the synthetic
background rate density.

The animation does not compare one smooth theoretical curve with another. It
repeats the actual detector prediction used by the objective.

For a sampled subset of the population, the same calculation produces either
individual population curves or a 10th--90th percentile envelope. That envelope
is deliberately described as a sampled population envelope.

It is not a credible interval.

The GA population is not a posterior distribution, no matter how attractively
Matplotlib shades it.

\section{How to read the replay}

The saved replay contains four connected views.

\subsection*{The folded spectrum}

The main panel shows the synthetic Poisson count-rate density, the current best
response-folded model and the sampled population envelope. This is the most
direct view of the phenotype produced by the current chromosome. Purple open
circles show the real MAXI/GSC reference on its native 2--10 keV grid. They are
context, not extra members of the fitness calculation.

\subsection*{The residuals}

The second panel uses Pearson-style residuals,

\begin{equation}
  r_i = \frac{n_i-\mu_i}{\sqrt{\max(\mu_i,1)}}.
\end{equation}

Channels outside the 0.5--10 keV objective band remain visible and are shaded.
They are not included in the C-statistic, but removing them from the figure
would make the detector appear to end exactly where the analysis becomes
convenient. The MAXI/GSC points do not enter this panel either: a raw difference
between two instruments is not a residual.

\subsection*{The score history}

The third panel follows both the best and median C-statistic. A champion can
improve while the rest of the population remains scattered across several bad
regions. Watching the median makes that harder to miss.

\subsection*{The decoded parameters}

The final panel reports the current values of $T_{\mathrm{in}}$,
$N_{\mathrm{disk}}$, $\Gamma$ and $K$. A falling score is useful. Seeing which
physical trade-off produced it is more useful.

The frame is therefore not four unrelated diagnostic plots sharing a timestamp.
Every panel refers to the same evaluated generation.

This is a surprisingly easy invariant to lose once plotting code becomes
enthusiastic.

\section{What did the E08 smoke run actually show?}

The demonstration run uses the smoke profile: 32 individuals and 24 evolved
generations. The raw GA consumes 800 objective evaluations, including the
initial population; the L-BFGS-B polish adds another 222. The replay contains
25 frames because generation zero is a scientifically meaningful state rather
than an inconvenient preface.

The headline results are shown in Table~\ref{tab:e08}.

\begin{table}[htbp]
  \centering
  \caption{Injected and fitted parameters for the E08 synthetic smoke run.}
  \label{tab:e08}
  \small
  \begin{tabular}{lrrrrr}
    \toprule
    Stage & $T_{\mathrm{in}}$ & $N_{\mathrm{disk}}$ & $\Gamma$ & $K$ & C-stat \\
    \midrule
    Synthetic truth       & 0.650000 & 8000.00 & 2.15000 & 0.270000 & -- \\
    Raw 24-generation GA  & 0.787671 & 1683.73 & 2.36647 & 0.388871 & 2079.873 \\
    GA + SciPy L-BFGS-B   & 0.648885 & 8060.95 & 2.15046 & 0.269375 & 180.495 \\
    \bottomrule
  \end{tabular}
\end{table}

The raw GA did not converge. It reached the smoke profile's maximum generation
with C-stat 2079.873.

That is not a minor implementation detail hidden in a JSON file.

It is the result.

The population found a useful region, but the champion had not resolved the
strong trade-off between disk temperature, disk normalization, photon index and
continuum normalization. The animation makes this much more obvious than a
single final parameter row.

Bounded SciPy L-BFGS-B then started from that champion and reached C-stat
180.495. The retained polished parameters lie extremely close to the injected
values.

This does not prove that the GA defeated SciPy.

It demonstrates a hybrid division of labour. The population searches globally
for a useful basin. A local bounded method finishes the numerical sentence more
efficiently than asking mutation to place the final decimal points.

\ifx\EvoComparisonPath\empty
  \begin{scopebox}[Figure unavailable]
  Run \texttt{python scripts/build\_blog\_bundle.py} before compiling this
  document to place the generated comparison image in \texttt{../assets/}.
  \end{scopebox}
\else
  \begin{figure}[htbp]
    \centering
    \includegraphics[width=\textwidth]{\EvoComparisonPath}
    \caption{Static E08 comparison. The upper panel overlays the synthetic
    target, the raw GA model, the GA + SciPy model and the real MAXI/GSC
    reference as purple open circles. The lower panel shows Pearson residuals
    for the synthetic target and models only, and shades detector channels
    outside the objective fit band. The real points are not included in the
    reported C-statistic or residuals.}
    \label{fig:comparison}
  \end{figure}
\fi

\section{Live calculation and saved replay are different jobs}

The live viewer is attached directly to the generation callback. On an
interactive Matplotlib backend it updates the current spectrum and residuals
while the optimiser is still calculating.

That answers a very immediate question:

\begin{quote}
What is the optimiser doing now?
\end{quote}

The saved replay is built from the completed population history. It works in a
headless environment and can be exported as HTML, GIF through Pillow, or MP4
when FFmpeg is available.

That answers a different question:

\begin{quote}
What exactly did this run do?
\end{quote}

The two paths use the same model evaluation, but they have different failure
modes. A live GUI depends on an interactive desktop backend. A saved file needs
to preserve every frame, avoid silently exceeding an embed limit and remain
useful after the Python process has ended.

The HTML writer produces a convenient self-contained Matplotlib fragment. For
the copy-ready website bundle, a small compiler extracts its 25 embedded PNG
frames, writes them as ordinary assets and replaces the inline player with
external CSS and JavaScript.

This removes the base64 overhead, permits a strict content-security policy and
adds a labelled timeline, keyboard-operable buttons, a paused initial state and
a static fallback.

Scientific Python is very good at making figures.

It is not automatically good at making those figures behave politely inside
somebody else's website.

\section{Where does real data enter?}

This is the section where it would be very easy to overstate what changed.

The replay includes a real MAXI/GSC spectrum of MAXI J1820+070 produced by the
official RIKEN/JAXA MAXI on-demand service. It covers MJD
58301.5--58302.5, close to the synthetic E08 anchor at MJD 58302.1, with
1539.5 seconds of exposure. The response EBOUNDS channels were grouped into 16
half-keV display bins across 2--10 keV.

That band contains 3004 source-region counts and 330 background-region counts.
The background is scaled using the source and background exposures and
\texttt{BACKSCAL} values before subtraction. The plotted errors propagate the
source and scaled-background Poisson uncertainties.

The derived CSV, exact request and provenance record travel with the website
bundle. The official result, including the source spectrum, background spectrum
and response, is available at
\url{https://maxi.riken.jp/mxondem/api/results/20260828013000_evoxrbv020real01/index.html}.

Version 0.2.0 can also load another observational spectrum from CSV and draw it
on its native energy grid beside the evolving synthetic fit. The preferred
schema is:

The preferred schema is:

\begin{lstlisting}
energy_keV,count_rate_density,count_rate_error
0.55,42.7,1.3
0.60,47.1,1.4
...
\end{lstlisting}

A counts table can instead provide \texttt{counts}, \texttt{exposure\_s} and
either \texttt{bin\_width\_keV} or lower and upper energy edges. Counts are
converted to count-rate density for plotting.

But the real spectrum remains a visual overlay only. The synthetic target is
what the GA and bounded SciPy L-BFGS-B fit. The MAXI values never change
selection and never enter the objective or the residual calculation.

\begin{warningbox}
\textbf{Two count-rate spectra on one axis are not automatically the same
measurement.}

A physical cross-instrument comparison needs the matching source and background
products, exposures, responses, effective areas, calibration choices and
likelihood treatment for each instrument. The vertical separation between the
MAXI/GSC points and the NICER-inspired synthetic model is therefore not a model
residual. EvoXRB deliberately refuses to turn that raw overlay into a fit
statistic.
\end{warningbox}

The reference makes the comparison real. The boundary around what was fitted
makes it accurate. The purpose of the loader is to preserve the distinction
between ``this curve can be drawn'' and ``this dataset has been calibrated well
enough to fit.''

Those are not the same claim.

\textbf{Acknowledgment:} This research has made use of MAXI data provided by
RIKEN, JAXA and the MAXI team.

\section{Checkpointing without changing the experiment}

A long genetic-algorithm run should be resumable.

It should not be resumable against different data while pretending that nothing
changed.

The animated workflow therefore binds each checkpoint to an objective signature
covering:

\begin{itemize}
  \item the case-study configuration and continuum choice;
  \item the target counts, exposure and fit mask;
  \item the true and detector energy grids and bin edges;
  \item effective area, redistribution matrix and background.
\end{itemize}

A signed checkpoint can only resume when the same signature is supplied. Old
generic checkpoints without signatures remain usable by generic callers, but a
domain workflow does not get to omit the signature after one has been recorded.

This is intentionally conservative.

Reusing a population is easy. Reusing cached objective values after changing
the response is wrong.

The run summary records both the short configuration digest and the complete
objective signature. The website compiler then emits selected public metadata
without absolute workstation paths or checkpoint locations.

Reproducibility is not improved by publishing the name of my Windows user
directory.

\section{Running and publishing it}

From the repository root, the normal smoke replay is generated with:

\begin{lstlisting}[language=bash]
python -m pip install -e ".[test]"
python -m evoxrb animate --profile smoke --epoch E08
\end{lstlisting}

The calculation can be watched live with:

\begin{lstlisting}[language=bash]
python -m evoxrb animate --profile smoke --epoch E08 --live
\end{lstlisting}

A user reference can be added with:

\begin{lstlisting}[language=bash]
python -m evoxrb animate \
  --profile smoke \
  --epoch E08 \
  --reference-csv data/reference/maxi_j1820p070_mjd58302.csv
\end{lstlisting}

The website bundle is compiled from the generated replay with:

\begin{lstlisting}[language=bash]
python -m evoxrb animate --profile smoke --epoch E08 --reference-csv data/reference/maxi_j1820p070_mjd58302.csv --output results/animations/E08_ga_spectra.html --comparison-output results/animations/E08_ga_comparison.png
python scripts/build_blog_bundle.py
python -m http.server 8000 --directory website/evoxrb-v0.2.0
\end{lstlisting}

The output uses relative paths and no remote runtime dependencies. Copying the
entire \texttt{website/evoxrb-v0.2.0} directory into a static website repository
moves the article, accessible replay, 25 frames, comparison plot, sanitized run
metadata and this LaTeX source together.

The source can be compiled with a conventional portable pdfLaTeX toolchain:

\begin{lstlisting}[language=bash]
cd website/evoxrb-v0.2.0/article
latexmk -pdf evoxrb-v0.2.0.tex
\end{lstlisting}

Running \texttt{pdflatex} twice is also sufficient when \texttt{latexmk} is not
available.

\section{So, did adding an animation improve the science?}

Did it make the raw smoke-profile GA converge?

\textbf{No.}

Did it make the algorithm look more competent than the scores say it was?

\textbf{Also no.} The structured residuals are now much harder to ignore.

Did it make the relationship between population search, parameter degeneracy
and local polishing easier to understand?

\textbf{Yes.}

Did it actually place a real reference spectrum beside the synthetic experiment
without quietly pretending that the educational response calibrated it?

\textbf{Yes.}

The next important step is still the one identified in Part I: increase the
difficulty of the likelihood landscape, then eventually replace the synthetic
products with a properly reduced observation and its matching calibration
files.

A four-parameter educational model remains a friendly problem by astronomical
standards. The stronger test will involve more physical Comptonization and
reflection components, background and calibration uncertainty, cross-instrument
constants, additional nuisance parameters and the multimodality that tends to
arrive once a model becomes interesting.

But v0.2.0 makes the intermediate experiment observable.

That matters because optimisation is not only a machine for producing a final
row of parameters.

It is a process.

And processes are much easier to trust when they are allowed to leave evidence.

\section{References and further reading}

\begin{enumerate}[label={[\arabic*]}]
  \item S. Seifi, \textit{EvoXRB: Genetic Algorithms, and a Very Real Black
  Hole} (2026). Part I:
  \url{https://svsam.com/blog/evoxrb-genetic-algorithms-x-ray-astronomy/}.

  \item W. Cash, \textit{Parameter estimation in astronomy through application
  of the likelihood ratio}, Astrophysical Journal 228, 939 (1979).
  \url{https://ntrs.nasa.gov/citations/19790043693}.

  \item SciPy documentation, \textit{minimize(method='L-BFGS-B')}.
  \url{https://docs.scipy.org/doc/scipy/reference/optimize.minimize-lbfgsb.html}.

  \item Matplotlib documentation, \textit{Animation API and writer classes}.
  \url{https://matplotlib.org/stable/api/animation_api.html}.

  \item NASA HEASARC, \textit{NICER analysis documentation}.
  \url{https://heasarc.gsfc.nasa.gov/docs/nicer/}.

  \item M. Matsuoka et al., \textit{The MAXI Mission on the ISS: Science and
  Instruments for Monitoring All-Sky X-Ray Images}, Publications of the
  Astronomical Society of Japan 61, 999 (2009).
  \url{https://ui.adsabs.harvard.edu/abs/2009PASJ...61..999M/abstract}.

  \item RIKEN, JAXA and the MAXI team, \textit{MAXI on-demand processing:
  MAXI J1820+070, MJD 58301.5--58302.5}.
  \url{https://maxi.riken.jp/mxondem/api/results/20260828013000_evoxrbv020real01/index.html}.
\end{enumerate}

\vfill
\begin{center}
\small\textcolor{EvoMuted}{
EvoXRB v0.2.0 --- synthetic / NICER-inspired --- Part II
}
\end{center}

\end{document}
